\documentclass{article}
\title{MAT1105 ukesoppgaver KMG}
\author{borgebj}
\date{}

\usepackage[norsk]{babel}
\usepackage{subcaption}
\usepackage{mathtools}
\usepackage{amssymb}
\usepackage{graphicx}
\usepackage[a4paper, top=1in, bottom=1in, left=1in, right=1in]{geometry}
\usepackage{float}
\usepackage{enumerate}



\setlength{\parskip}{2em}
\setlength{\parindent}{0pt}


\begin{document}

\maketitle

\section*{Kapittel 2}

\subsection*{Oppgave 2.1}

(a, b) når $b < a$ er: 
\[ (a, b) = \{ x \in \mathbb{R} \mid x > a\; og \; x < b\} \]

\subsection*{Oppgave 2.2}

Mengden av alle positive partall er: 

1) \textbf{Liste:} \hspace{6em} 
\(P = \{2, 4, 6, 8, ...\}\) 

2) \textbf{Mengdebygger:} \hspace{1em}
\( P = \{ 2x \mid x \in \mathbb{N}\} \) \\

\subsection*{Oppgave 2.3}

\(A = \{x \in \mathbb{R} \mid x > 4\} \hspace{2em} \text{ og } \hspace{2em} B = \{x \in \mathbb{R} \mid x \ge 0\}\)

Vi velger et vilkårlig tall 'x' fra A.\\
Dette tallet 'x' er høyere enn 4.\\
Ethvert tall i B er høyere 0.\\
'x' er høyere enn 0 siden 'x' er høyere enn 4, derfor er 'x' en delmengde av B\\

\newpage

\subsection*{Oppgave 2.4}

A = B er det samme som \(A \subseteq B \ og \ B \subseteq A\)

Dette er fordi hvis A er lik B så må alt som er i A være med i  B\\
Hvis B er lik A så må alt som er i B være med i A også.\\

\subsection*{Oppgave 2.5}

\(A = \{1, 2, 3\}, \  B = \{3, 1, 1, 2\}\)

\(
1) \ A \subseteq B \\ \\
1 \in B ? \ Ja \\
2 \in B ? \ Ja \\
3 \in B ? \ Ja \\ \\
2) \ B \subseteq A \\ \\
3 \in A ? \ Ja \\
1 \in A ? \ Ja \\
1 \in A ? \ Ja \\
2 \in A ? \ ja 
\)

Slikt er det vist at alle elementer i A er i B, og er derfor en delmengde av B.

\subsection*{Oppgave 2.6}
\(A = \{1, 2\} \ og \ B = \{1, 2, \{1, 2\}\}\)

\begin{enumerate}[a.]
\item \(A \subseteq B\) \ - \textbf{Sant} \ \(1 \in B \text{ og } 2 \in B\)
\item \(A \in B\) \ - \textbf{Sant} \ \(\ A = \{1, 2\} \ og \ \{1,2\} \in B\)
\item \(B \subseteq A\) \ - \textbf{Ikke sant} \ \( \{1, 2\} \notin A\)
\item \(B \in A\) \ - \textbf{Ikke sant} \ \(\{1, 2, \{1, 2\}\} \notin A\)
\end{enumerate}

\newpage

\subsection*{Oppgave 2.7}

\begin{enumerate}[a.]
    \item \(A = \emptyset, \qquad\qquad\qquad\quad \underline{|A| = 0}\)
    \item \(B = \{8, 10, 3, 5\}, \quad\quad\ \ \underline{|B| = 4}\)
    \item \(C = \{2, 5, 2, 9, 5\}, \qquad \ \underline{|C| = 5}\)
    \item \(D = \{x \in \mathbb{N} \mid x \le 20 \text{ og x er et partall}\}, \quad \underline{|D| = 10}\)
    \item \(E = \{x \in \mathbb{Z} \mid x < 0\}, \quad \underline{|E| = \infty}\)
\end{enumerate}

\subsection*{Oppgave 2.8}

\begin{enumerate}[-]
    \item La a, b, c, d være fire tall.
    \item La \(A = \{a, b, c, d\}\) 
    \item \#A forteller oss om kardinaliteten, aka 'størrelsen' 
    \item I dette tilfellet er \#A = 4 
\end{enumerate}

\subsection*{Oppgave 2.9}
\[A = \{x \in \mathbb{R} \mid x > 4 \}, \quad B = \{x \in \mathbb{R} \mid x < 9 \}\]

\begin{enumerate}
    \item \(A \cup B = \{x \in \mathbb{R} \mid x > 4 \text{ eller } x < 9\}\)
    \item \(A \cap B = \{5, 6, 7, 9\} = (4, 9)\)
    \item \(A / B = \{x \in \mathbb{R} \mid x \ge 9 \} = [9, \infty)\)
\end{enumerate}

\subsection*{Oppgave 2.10}
Hvis A og B er mengder:

\begin{enumerate}[a.]
    \item \(\mathbf{A \cap B = B \cap A}\)
    \begin{enumerate}[-]
        \item Sann fordi \(A \cap B\) er alt som er i både A og B.
        \item \(B \cap A\) sier det samme, det er alt som er i både B og A.
        \item Sann fordi snitt er kommutativ.
    \end{enumerate} 
    \item \(\mathbf{A \cup B = B \cup A}\)
    \begin{enumerate}[-]
        \item Stort sett samme som ovenfor, fordi:
        \item \(A \cup B\) er alt som er i både A og B i en ny mengde
        \item \(B \cup A)\) er alt som er i både B og A, det samme.
        \item Sann fordi union er kommutativ
    \end{enumerate} 
    \item \(\mathbf{A \cap B \subseteq A}\)
    \begin{enumerate}[-]
        \item \(A \cap B\) evaluerer til alt som er i både A og B.
        \item Om den ikke er tom, er det alltid en delmengde av A siden fra definisjon av snitt er den i både A og B.
        \item Om den er tom, den tomme mengden, er det også en delmengde av A, siden den tomme mengden er delmengde av alle mengder.
    \end{enumerate}
    \item \(\mathbf{A \ \backslash\ B \subseteq A}\)
    \begin{enumerate}[-]
        \item \(A \backslash B\) vil si at vi fjerner alt som er i B fra A. 
        \item Om ingen elementer fra B er i A, så fjernes ingenting og vi ender opp med samme mengde A.
        \item Om elementer blir fjernet, ligger vi fremdeles igjen med en delmengde av A, om vi ender opp med elementer eller den tomme mengden
    \end{enumerate}
    \item \(\mathbf{A \subseteq A \cup B}\)
    \begin{enumerate}[-]
        \item \(A \cup B\) gir oss en ny mengde som inneholder elementer fra \underline{både A og B}, derfor er det alltid en delmengde av A.
    \end{enumerate}
\end{enumerate}

\subsection*{Oppgave 2.11}
\((A^c)^c = A\) \\ \\
\(A^c\) forteller oss her at det er mengden av ALT som \underline{ikke} er med i mengden A. \\ \\
\((A^c)^c\) vil da si mengden av alt som ikke er med i mengden av alt som ikke er med i A. \\ \\
Da sier det seg selv at mengden må være A, fordi hvis mengden inneholder alt som ikke er inneholdt "utenfor" A, må den være "innenfor" A.

\subsection*{Oppgave 2.12}
\begin{enumerate}[a.]
    \item \(\mathbf{A = (3, \infty)}\)
    \begin{enumerate}[-]
        \item \(A^c = \mathbb{R} \backslash A = \{x \in \mathbb{R} : x \notin (3, \infty)\} = \{x \in \mathbb{R} : x \le 3\} = (-\infty, 3)\) 
    \end{enumerate}

    \item \(\mathbf{B = \{2, 4, 6, 8, ...\}}\)
    \begin{enumerate}[-]
        \item \(B^c = B \backslash \mathbb{N} = \{n : n \notin \{2, 4, 6, 7, 8, ... \}\} = \{2n+1 \: n \in \mathbb{N}\} = \{1, 3, 5, 7, ...\}\)
    \end{enumerate}

    \item \(\mathbf{C = \mathbb{Q}}\)
    \begin{enumerate}[-]
        \item $\mathbb{Q}$ er alle rasjonale tall
        \item Vi lar universet være $\mathbb{R}$
        \item \(C^c = \mathbb{R} \ \backslash \ \mathbb{Q}\) = alle irrasjonale tall.
    \end{enumerate}
\end{enumerate}

\newpage

\subsection*{Oppgave 2.13}

\((A \cup B)^c = A^c \cap B^c\) = "Komplementet av unionen mellom A og B er lik snittet mellom komplementene av A og B" \\

\((A \cap B)^c = A^c \cup B^c\) = "Komplementet av snittet mellom A og B er lik unionen mellom komplementene av A og B

\subsection*{Oppgave 2.14}

\(A = \{x \in \mathbb{R} : x < 0\} \text{  og  } B = \{x \in \mathbb{R} : x \ge 1\}\)

\((A \cup B)^c = A^c \cap B^c \\ \\
A^c = \{x \in \mathbb{R} : x \ge 0\} = [0, \infty)\\
B^c = \{x \in \mathbb{R} : x < 1\}  = (-\infty, 1] \\ \\
A^c \cap B^c = B^c \cap A^c = (-\infty, 1] \cap [0, \infty) = [0, 1]\)

\subsection*{Oppgave 2.15}

Forskjellen mellom n-tuppelet \(A = (x_1, x_2, ..., x_n) \text{ og mengden } B = \{x_1, x_2, ..., x_n\}\) er hovedsakelig 2 ting:
\begin{enumerate}
    \item Rekkefølgen på elementene. i n-tuppelet A så spiller rekkefølgen på elementene en rolle, og vi kan ikke endre den. I mengden B spiller ikke rekkefølge noen rolle og vi kan skrive den som vi vil, e.g. \(B = \{x_2, x_1, ...,  x_n\}\)
    \item Forekomster av elementer. I n-tuppelet A så har det noe å si hvis vi har 2 av samme element, mens i mengden B vil 2 av samme element ikke spille noen rolle, og bety kun at det elementet er i mengden. Kan forkortes til kun ett av det elementet.
\end{enumerate}

\subsection*{Oppgave 2.16}

\(A = \{1, 2, 3\} \text{ og } B = \{0, 4\}\)

\(A \times B = \{(1, 0), (1, 4), (2, 0), (2, 4), (3, 0), (3, 4)\}\)

\(\#A = 3\quad \#B = 2 \quad \#A(\times B) = 3 \cdot 2 = 6\)

\subsection*{Oppgave 2.17}

A og B er endelige mengder.

\(A \times B\) inneholder da \#A $\cdot$ \#B elementer

\section*{Kapittel 3}

\subsection*{Oppgave 3.1}

\(f : A \rightarrow B\) \\
A = definisjonsområdet \\
B = verdiområdet \\

\(V_f = \{ f(x) : x \in A \} \) = verdimengden

Fra definisjonen til en funksjon, er alltid f(x) med i B.
\\ Dermed er alle punkter i verdimengden også i B.
\\ Ethvert objekt i verdimengden er del av verdiområdet, siden ethvert objekt i funksjonen, f(x) er tatt ifra verdiområdet.

Eksempel: \\
\( A = \{a, b, c\} \quad B = \{1, 2, 3\} \\ f : A \rightarrow B = \{(a, 1), (b, 2), (c, 1)\}\)

Eksempel: \\
\(A = \mathbb{R} \quad B = \mathbb{R} \\ f : A \rightarrow B = f : \mathbb{R} \rightarrow \mathbb{R}\) \\ 
f.eks. funksjonen $f(x) = x^2$ \\
Verdiområdet er $\mathbb{R}$ \\
Verdimengden er \([0, \infty)\)

\subsection*{Oppgave 3.2}
Kardinalitetssymbolet \# 'mapper' alle mengder til et tall, som beskriver størrelsen. \\ \\
Definisjonsområdet til kardinalitetssymbolet \# er alle mengder. \\ \\
Verdiområdet til \# er mengden av alle heltall. \\ \\

\newpage

\subsection*{Oppgave 3.3}

\subsection*{Oppgave 3.4}
En funksjon $f : A \rightarrow B$ er bijektiv hvis det for enhver $y \in B$ finnes én, og bare én, $x \in A$ slik at $f(x) = y$

En funksjon er injektiv hvis $f(x_1) \neq f(x_2)$ for alle $x_1, x_2 \in A$ som er slik at $x_1 \neq x_2$
\\ - Dvs. at hver tilording / bruk av funksjonen gir et "unikt" element fra verdimengden. 
\\ - Man skal ikke komme fram til samme resultat ved bruk av 2 ulike argumenter.
\\ \\\
En funksjon er surjektiv for enhver $y \in B$ finnes en $x \in A$ slik at $f(x) = y$
\\ - Dvs. at enhver verdi i B, verdimengden, så finnes en verdi/et argument i A, definisjonsmengden.
\\ \\
Ved å si at for enhver $y \in B$ finnes én og bare én $x \in A$ slik at $f(x) = y$ beskriver vi både at enhver verdi i B skal ha en verdi i A som oppfyller surjektivitet, og at det finnes kun én og bare én tilordning av dette, oppfyller injektivitet.

\subsection*{Oppgave 3.5}
\(f : \mathbb{R} \rightarrow B \text{ gitt ved } f(x) = x^2\)
\\ - Ved å si at $B = [0, \infty)$ så gjør vi funksjonen surjektiv. 
\\ - Dette er fordi hadde vi hatt $B = \mathbb{R}$ kunne vi aldri fått negative tall. 
\\ - Med ingen negative, kan vi bruke alle relle tall til å få et resultat som er enten 0 eller positivt.
\\ - Den er fortsatt ikke injektiv, fordi vi kan fortsatt sette inn negative tall, fordi $A = \mathbb{R}$, og pga $x^2$ så resulterer et tilsvarende positivt og negativt tall det samme.

\subsection*{Oppgave 3.6}

Surjektivitet : En funksjon $f : A \rightarrow B$ er surjektiv hvis det for alle $y \in B$, finnes en $x \in A$ slik at $f(x) = y$
\\
Ved å si at verdiområdet B og verdimengden til funksjonen er like, sier vi at alle elementer i B skal være 'truffet', i andre ord ha en x slik at y er i verdimengden

\subsection*{Oppgave 3,7}
\newpage

\section*{Kapittel 4}

\subsection*{Oppgave 4.1}
Er følgende utsagn, predikat, eller ingen?

\begin{enumerate}[(a)]
    \item $2 \in \mathbb{N}$ \\ Dette er et utsagn, og den er sann.
    \item \say{Alle partall er delelig med 2} \\
    Dette er et utsagn, og er sant fordi det følger definisjonen av partall.
    \item \say{Det finnes liv på andre planeter} \\
    Dette er et utsagn, og er umulig å si sann eller usann.
    \item $3x + 8 = 17$ \\
    Dette er et predikat, på grunn av variabelen x, og vi kan derfor ikke si noe mer om det.
    \item \say{Mengdene $\mathbb{N}$ og $\mathbb{Z}$ er polymorfiale, men ikke $\mathbb{R}$} \\
    
\end{enumerate}

\end{document}
